\begin{table}[p]\centering
    \caption{\\Aggregated investment effects}\footnotesize
    \label{tab:agg}
    \begin{tabular}{lrrr}
    \toprule
    & 1933 & 1934 & After \\
    \midrule
    \multicolumn{4}{c}{\textit{Panel A: All firms}}\\
    \midrule
    Total net investment in \%               & -5.47 & -2.47 & -0.21 \\
    Gold clause effect in \%               & -1.79 & -1.23 & -0.15 \\
    \midrule
    \multicolumn{4}{c}{\textit{Panel B: $\tilde{d}>0$ firms}}\\
    \midrule
    Total net investment in \%               & -6.59 & -3.01 & -0.75 \\
    Gold clause effect in \%               & -3.27 & -2.18 & -0.28 \\
    \midrule
    \multicolumn{4}{c}{\textit{Panel C: $\tilde{d}=0$ firms}}\\
    \midrule
    Total net investment in \%               & -4.11 & -1.78 & 0.56 \\
    \bottomrule

    \end{tabular}\\

    \vspace*{3mm} \justifying \noindent
\footnotesize{\textit{Notes.} This table aggregates firm-level estimates to assess the total investment foregone due to gold clause uncertainty. Total net investment is the aggregate change in fixed capital divided by lagged aggregate fixed capital, for firms present in both years. The gold clause effect is the capital-weighted average of each firm's foregone investment, calculated using the regression coefficients from Table \ref{tab:inv_main} applied to each firm's gold clause exposure $\tilde{d}$ (as defined in equation (\ref{eq:tilde_d})). Panel A reports results for all firms; Panels B and C decompose by gold clause exposure. ``After'' denotes the post-resolution period (1935--1940). See Section \ref{sec:agg} for details.}
\end{table}